\section{Conclusion}

The aim of this project was to design, implement, and evaluate a functioning Micromouse robot capable of autonomous maze navigation within a reduced academic project scope. Within this scope, the aim was achieved. The final system combined a custom embedded hardware platform with closed-loop motion control, infrared wall sensing, encoder-based odometry, and map-based navigation software, and it was able to explore a reduced $6 \times 6$ maze and use the generated map for a later goal-directed run. As a result, the project demonstrated that the essential tasks of sensing, control, localisation, and navigation could be integrated successfully into a single autonomous robotic platform.

From a technical point of view, the main achievement of the project was not maximum speed, but the successful integration of multiple engineering disciplines into one coherent system. The implemented hardware architecture, based on a dsPIC-centered PCB with dedicated sensing, drive, communication, and power subsystems, proved suitable for this purpose. In addition, the layered software structure and the cascaded controller design provided an effective basis for organising the real-time behaviour of the robot and for managing the interaction between low-level actuation and high-level navigation.

At the same time, the project also revealed the limits of the current implementation. The robot should be understood as a functional academic prototype rather than as a fully optimised competition Micromouse. The most important technical limitations were the approximate open-loop turning behaviour, the sensitivity of cell-centre detection and localisation to motion disturbances, and the deliberately reduced maze scope. Several practical issues during development, including the programming-interface error, Bluetooth integration problems, and limited wireless range, also increased the implementation effort and reduced debugging convenience. Taken together, these points show that the implemented concept was technically viable, but that further refinement would be required for a more demanding and general Micromouse system.

\subsection{Outlook}

Several directions for future work follow directly from the present results. A first important step would be the extension of the navigation software from the current reduced maze representation to the full classical $16 \times 16$ Micromouse problem. In parallel, the motion performance could be improved through dedicated closed-loop turning control, more robust localisation, and continued refinement of the wall-based trajectory correction. On the hardware side, a revised programming interface and improved wireless communication layout would increase maintainability and debugging convenience. Taken together, these improvements would build naturally on the implemented prototype and could transform it from a successful educational platform into a more capable and competition-oriented autonomous robot.
